\section{Implementation}

\label{sec:implementation}
% ══════════════════════════════════════════════════════════════════════════════

\subsection{Technology Stack}

\begin{itemize}
  \item \textbf{Mobile apps}: React Native (iOS/Android) with offline
    observation caching.
  \item \textbf{Backend}: FastAPI (Python) with PostgreSQL (PostGIS) for
    spatial data.
  \item \textbf{AI inference}: TaxoNet served via ONNX Runtime on
    NVIDIA T4 GPUs; 45\,ms per image.
  \item \textbf{Gamification engine}: PyTorch-based contextual bandit
    model; daily retraining.
  \item \textbf{Quality model}: Bayesian inference using NumPyro;
    batch updates every 6 hours.
  \item \textbf{Blockchain}: ZooRep token distribution via Zoo Network
    smart contracts.
  \item \textbf{Data storage}: Observations in PostgreSQL; images in
    S3-compatible object storage; IPFS for immutable research-grade
    data snapshots.
\end{itemize}

\subsection{Scalability}

The platform processes approximately 10{,}000 observations per hour at
peak usage. AI inference is the bottleneck, requiring 3 T4 GPUs at peak.
The quality model runs asynchronously and does not affect submission
latency. Challenge generation is precomputed in nightly batch runs and
cached for real-time delivery.


% ══════════════════════════════════════════════════════════════════════════════
